\documentclass[10pt,a4paper]{amsart}
\usepackage[margin=19mm]{geometry}
\usepackage{amsmath,amssymb,mathtools}
\usepackage[scr=rsfs,cal=euler]{mathalfa}
\usepackage{xfrac}
\usepackage[unicode,hidelinks,bookmarks=false]{hyperref}
\newcommand{\ie}{i.e.\ }
\newcommand{\cf}{cf.\ }
\newcommand{\resp}{respectively\ }
\newcommand{\Subsection}{\subsection}
\providecommand{\nobreakdash}{\nobreak\hbox{-}}
\setlength{\emergencystretch}{2em}
\multlinegap0pt
\usepackage{bm}
\usepackage{bbm}
\usepackage{dsfont}
\usepackage{xspace}
\usepackage{cancel}
\usepackage{mathtools}
\usepackage{xcolor}
\usepackage{amsthm}
\makeatletter
\@ifundefined{theorem}{\newtheorem{theorem}{Theorem}}{}
\@ifundefined{assumption}{\newtheorem{assumption}[theorem]{Assumption}}{}
\makeatother
\makeatletter
\expandafter\ifx\csname customass\endcsname\relax
\newenvironment{customass}[1]
{\renewcommand\theinnercustomass{#1}\innercustomass}
{\endinnercustomass}
\fi
\expandafter\ifx\csname customthm\endcsname\relax
\newenvironment{customthm}[1]
{\renewcommand\theinnercustomthm{#1}\innercustomthm}
{\endinnercustomthm}
\fi
\newcommand*{\dd}{\mathrm{d}}
\newcommand\norm[1]{\left\lVert#1\right\rVert}
\newcommand\normm[1]{\big\lVert#1\big\lVert}
\makeatother
\title[The Influence of the Memory Capacity of Neural DDEs on the U]{The Influence of the Memory Capacity of Neural DDEs on the Universal Approximation Property}
\author{Christian Kuehn, Sara-Viola Kuntz}
\date{Source: arXiv:2505.07244v2 (CC BY 4.0)}
\begin{document}
\maketitle

\noindent\emph{Source and licence.} The Influence of the Memory Capacity of Neural DDEs on the Universal Approximation Property. Version arXiv:2505.07244v2 (CC BY 4.0), distributed under the Creative Commons Attribution 4.0 International licence, as recorded in the arXiv licence metadata for this exact version and shown on the abstract page https://arxiv.org/abs/2505.07244v2. Full rights notice: licenses/arxiv-2505.07244-CC-BY-4.0.txt.\par\smallskip

\noindent\textit{Editorial scope.} Counted unit: the Theorem whose source label is \texttt{th:universalaugmented} in the article (source0.tex lines 1046--1053, source label \texttt{th:universalaugmented}), delivered together with the article's own complete original proof of it (source0.tex lines 1055--1123, one \texttt{proof} environment of 69 lines). No mathematics is written, repaired or re-derived: statement and proof are the article's own text line for line. Required context retained uncounted: eq:DDEconstantIC (lines 617--622); eq:DDEaffinelinear (lines 628--630); ass:vectorfield (lines 634--636). Every definition, display and label of those lines is retained unchanged. Omitted from this package and not redistributed: 1--616, 623--627, 631--633, 637--1045, 1054--1054, 1124--2419. Nothing inside a retained excerpt is omitted or altered: every retained body is delivered exactly as the article's lines are written, so no omission receipt of any kind accompanies this package. Citations: the retained excerpts carry no citation command at all, so no bibliography entry of the article is transcribed and no reference is redistributed. Reference adaptation: none was needed and none was made; every reference inside the delivered unit resolves inside it, so no unresolved reference or citation is produced. Printed slips kept literal: no printed word, symbol, hypothesis or number of the counted statement or of its proof was altered. Layout and class: the article's own class is \texttt{article} with packages including inputenc, amsmath, amsfonts, amssymb, amsthm, geometry, color, transparent, xcolor, graphicx, subcaption, tabularx, float, hyperref; the wrapper is the lane's pinned \texttt{amsart} editorial wrapper at 10pt on A4, which restates the theorem-like environments the retained text needs and adds the article's own macro definitions that the retained text needs (\texttt{\textbackslash dd}, \texttt{\textbackslash norm}, \texttt{\textbackslash normm}), with extra wrapper packages (bm, bbm, dsfont, xspace, cancel, mathtools, xcolor). The delivered numbering is the unit's own. Assets: no retained span contains a figure environment, a graphics-inclusion command, a PDF-page inclusion, a table or a TikZ picture; no image file is copied into this package, the package declares no asset dependency and no image file is redistributed with it. Licence: version arXiv:2505.07244v2 of this article is distributed under the Creative Commons Attribution 4.0 International licence, as recorded in the arXiv licence metadata for this exact version and shown on the abstract page https://arxiv.org/abs/2505.07244v2. A full rights notice is delivered at licenses/arxiv-2505.07244-CC-BY-4.0.txt. Characters: every retained excerpt is pure ASCII, so the delivered engine, XeLaTeX with its default Latin Modern text font, reports no missing character.
	\begin{equation} \label{eq:DDEconstantIC}
		\begin{aligned}
			\frac{\dd y}{\dd t} &= F(t,y_t) \qquad &&\text{for } t\geq 0, \\  
			y(t) &= c_{y_0}(t) &&\text{for } t \in [-\tau,0],
		\end{aligned}
	\end{equation}

	\begin{equation}\label{eq:DDEaffinelinear}
		\Phi: \mathcal{X} \rightarrow \mathbb{R}^q, \quad x \mapsto \tilde{\lambda}(y(0,c_{\lambda(x)})(T)) = \widetilde{W} \cdot y(0,c_{Wx+b})(T) + \tilde{b} ,
	\end{equation}

	\begin{assumption}[Neural DDE Well-Definedness]\label{ass:vectorfield}
		For the neural DDE architecture \eqref{eq:DDEaffinelinear} based on the DDE \eqref{eq:DDEconstantIC} it holds that $F:\Omega\rightarrow\mathbb{R}^m$ is defined on $\Omega = \Omega_t \times \Omega_y \subset \mathbb{R}\times \mathcal{C}$ open, and $[0,T]\subset \Omega_t $, $[-\tau,T]\subset \mathcal{I}_{y_0}$ for all $y_0\in \mathcal{A} = \lambda(\mathcal{X})$ and $\Omega_0 \coloneqq \left\{c_{y_0}\in \mathcal{C}: y_0 \in \mathcal{A}\right\} \subset \Omega_y$.
	\end{assumption}

	\begin{theorem}[Universal Embedding for Augmented Neural DDEs] \label{th:universalaugmented}
		Let $k \geq 0$ and $\Psi\in C^0(\mathcal{X}, \mathbb{R}^q)$ if $k = 0$ or $\Psi\in C^k_b(\mathcal{X},\mathbb{R}^q)$ if $k\geq 1$, with $\mathcal{X}\subset \mathbb{R}^n$ open and global Lipschitz constant $K_\Psi\in[0,\infty)$. Fix constants $m \geq n+q$, $\tau \in [0,T]$, $w,\tilde w\in (0,\infty)$ and $K \in [0,\infty)$, such that $$K T \geq \frac{K_\Psi}{w \tilde{w}}.$$ Then there exists an augmented neural DDE architecture $\Phi \in \textup{NDDE}_{\tau,\textup{A},K}^k(\mathcal{X},\mathbb{R}^q)$  embedding the map~$\Psi$ with the following properties:
		\begin{itemize}
			\item The underlying vector field $F: \mathbb{R}\times \Omega_y\rightarrow \mathbb{R}^m$  is globally Lipschitz continuous in the second component on $\mathbb{R}\times \Omega_0$  with Lipschitz constant $K\in[0,\infty)$, where $\Omega_0 = \{c_{\lambda(x)}\in\mathcal{C}: x\in\mathcal{X}\}\subset \Omega_y$ is the set of constant initial data used.
			\item The neural DDE $\Phi$ has weight matrices $W$, $\widetilde{W}$ with $\normm{W}_\infty = w$ and $\normm{\widetilde{W}}_\infty = \tilde w$.
		\end{itemize}
		Consequently, any globally Lipschitz continuous map $\Psi\in C^0(\mathcal{X},\mathbb{R}^q)$ or $\Psi\in  C^k_b(\mathcal{X},\mathbb{R}^q)$ can be embedded for every $\tau \in [0,T]$ in the augmented neural DDE architecture $\textup{NDDE}^k_{\tau,\textup{A},K}(\mathcal{X},\mathbb{R}^q)$with $m \geq n+q$.
	\end{theorem}

	\begin{proof}
		Fix $m \geq n+q$, $\tau\in[0,T]$, $w, \tilde w \in (0,\infty)$, $K \geq \frac{K_\Psi}{w \tilde w T}$, and define an augmented neural DDE $\Phi \in \textup{NDDE}_{\tau,\textup{A},K}^k(\mathcal{X},\mathbb{R}^q)$ with weights and biases $W = w \cdot  \textup{Id}_{m,n}\in\mathbb{R}^{m \times n}$, $b = 0\in\mathbb{R}^m$, $\tilde{b} = 0 \in \mathbb{R}^q$ and
		\begin{equation*}
			\widetilde{W} = \begin{pmatrix}
				0_{q,n} && \hspace{-4mm}\tilde w \cdot \textup{Id}_{q,q} && \hspace{-4mm} 0_{q,m-n-q}
			\end{pmatrix} \in \mathbb{R}^{q \times m}.
		\end{equation*}
		In order to embed the map $\Psi$, we define the DDE
		\begin{equation}\label{eq:DDE_universal_aug}
			\begin{aligned}
				\frac{\dd y}{\dd t} &= F(t,y_t) \qquad &&\text{for } t\geq 0, \\  
				y(t) &= c_{\lambda(x)}(t) &&\text{for } t \in [-\tau,0],
			\end{aligned}
		\end{equation}
		with constant initial data $c_{\lambda(x)}:[-\tau,0]\rightarrow \Omega_y$, $\lambda(x) = Wx$, $x \in \mathcal{X}$ and continuous vector field $	F: \mathbb{R}\times \Omega_y\rightarrow \mathbb{R}^m$ with 
		\begin{equation*}
			\Omega_y =   \left\{u\in\mathcal{C}: \tfrac{1}{w}\cdot u_{1,\ldots,n}(t) \in \mathcal{X} \text{ for } t \in [-\tau,0]\right\}.
		\end{equation*}
		$\Omega_y$ is an open subset of $\mathcal{C}$, as the first $n$ components of $u(t)$ are contained in the open set $w \cdot \mathcal{X}$ and the components $u_{n+1,\ldots,m}(t)$ are contained in the open set $\mathbb{R}^{m-n}$.  Consequently, the domain of definition  $\mathbb{R}\times \Omega_y$ of the vector field $F$ is an open subset of $\mathbb{R} \times \mathcal{C}$ and $\Omega_y$ contains the  set of initial functions used, given by $\Omega_0 = \{c_{\lambda(x)}\in\mathcal{C}:x\in\mathcal{X}\}$. The vector field $F$ is defined as 
		\begin{equation} \label{eq:vectorfield_aug}
			F(t,y_t) =\begin{pmatrix}
				\big[\frac{\dd y}{\dd t}\big]_{1,\ldots,n} \textcolor{white}{\ldots ....} \\
				\big[\frac{\dd y}{\dd t}	\big]_{n+1,\ldots,n+q}\textcolor{white}{.} \\
				\big[\frac{\dd y}{\dd t}\big]_{n+q+1,\ldots,m}
			\end{pmatrix}
			= 
			\begin{pmatrix}
				0 \textcolor{white}{\big[\frac{1}{1}\big]_1}\\
				\frac{1}{\tilde w T} \cdot \Psi\left(\frac{1}{w}\cdot y_{1,..,n}(t)\right)\\
				0 \textcolor{white}{\big[\frac{1}{1}\big]_1}
			\end{pmatrix}, \qquad t\in\mathbb{R},
		\end{equation}
		which is the vector field of an ODE independently of the delay $\tau\in[0,T]$. The initial data is only used at the point $t = 0$ and given by
		\begin{equation*}
			\begin{pmatrix}
				y_{1,\ldots,n}(0) \\
				y_{n+1,\ldots,n+q}(0) \\
				y_{n+q+1,\ldots,m}(0)
			\end{pmatrix}
			= \lambda(x) = Wx+b = \begin{pmatrix}
				wx \\ 0 \\ 0
			\end{pmatrix}.
		\end{equation*} 
		With the constant initial function $c_{\lambda(x)}$ on the time interval $[-\tau,0]$, the specified initial value problem can be interpreted as a DDE. For $t\in[0,T]$ and $x\in\mathcal{X}$ it holds $y_{1,\ldots,n}(t) = wx$, such that the vector field $F$ is well defined. As $\Psi\in C^0(\mathcal{X}, \mathbb{R}^q)$ if $k = 0$ or $\Psi\in C^k_b(\mathcal{X},\mathbb{R}^q)$ if $k\geq 1$, it follows that $F \in C^{0,0}(\Omega,\mathbb{R}^m)$ if $k = 0$ or $F\in C_b^{0,k}(\Omega,\mathbb{R}^m)$ if $k\geq 1$. The unique, continuous solution of \eqref{eq:DDE_universal_aug} is given by
		\begin{equation*}
			\begin{pmatrix}
				y_{1,\ldots,n}(t) \\
				y_{n+1,\ldots,n+q}(t) \\
				y_{n+q+1,\ldots,m}(t)
			\end{pmatrix}
			=  \begin{pmatrix}
				wx \\ \frac{t}{\tilde w T} \cdot \Psi(x) \\ 0
			\end{pmatrix},
		\end{equation*} 
		such that it follows with the affine linear layer $\tilde \lambda: y \mapsto \widetilde{W}y+b$ that
		\begin{equation*}
			\Phi(x) = \tilde{\lambda}(y(T))  = \widetilde{W}\cdot \begin{pmatrix}
				x \\ \frac{1}{\tilde w}\cdot \Psi(x) \\ 0
			\end{pmatrix} = \Psi(x) \qquad \text{for all } x \in \mathcal{X}.
		\end{equation*}
		Since for every $x \in \mathbb{R}^n$, we determined via explicit integration the unique and continuous solution of the corresponding DDE, which satisfies Assumption \ref{ass:vectorfield}, the defined neural DDE $\Phi$ embeds the map $\Psi$ and is an element of $\textup{NDDE}^k_{\tau,\textup{A},K}(\mathcal{X},\mathbb{R}^q)$. $\Phi$ has by construction a delay $\tau\in[0,T]$ and the defined vector field \eqref{eq:vectorfield_aug} is globally Lipschitz continuous in the second variable on $\mathbb{R}\times \Omega_y$ with Lipschitz constant $K$, as for all $t \in \mathbb{R}$ and $y_t,z_t \in \Omega_y$ it follows
		\begin{align*}
			\norm{F(t,y_t)-F(t,z_t)}_\infty  & \leq \norm{\frac{1}{\tilde w T}\Psi\left(\frac{1}{w}\cdot y_{1,..,n}(t)\right)- \frac{1}{\tilde w T}\Psi\left(\frac{1}{w}\cdot z_{1,..,n}(t)\right)}_\infty \\
			& = \frac{1}{\tilde wT}\norm{\Psi\left(\frac{1}{w}\cdot y_{1,..,n}(t)\right)- \Psi\left(\frac{1}{w}\cdot z_{1,..,n}(t)\right)}_\infty \\
			& \leq \frac{K_\Psi}{w \tilde w T} \norm{y_{1,..,n}(t)-z_{1,..,n}(t)}_\infty  \\
			& \leq K \norm{y(t)-z(t)}_\infty \leq K\norm{y_t-z_t}_\infty.
		\end{align*}
		In the proof we used in the second line that the map $\Psi$ is globally Lipschitz continuous on $\mathcal{X}$ with Lipschitz constant $K_\Psi\in[0,\infty)$ and the last line follows from the assumption that $\frac{K_\Psi}{w \tilde w T} \leq K$. 
	\end{proof}

\end{document}
